\chapter{Behavioural Economics}
\label{chap:behavioural}

\section{Learning objectives}

After completing this chapter, you should be able to:
\begin{itemize}
  \item explain how behavioural economics extends, rather than simply rejects, rational-choice analysis;
  \item explain the System 1--System 2 teaching framework and when deliberate checking is useful;
  \item define bounded rationality, heuristics, anchoring, availability, representativeness, and regression to the mean;
  \item explain reference dependence, framing, loss aversion, and present bias;
  \item distinguish overconfidence, optimism, planning fallacy, and status quo bias;
  \item read and interpret a simple reference-dependent value function;
  \item distinguish a nudge from a tax, a ban, and a standard information policy; and
  \item evaluate behavioural policy using effectiveness, autonomy, distribution, and legitimacy.
\end{itemize}

\section{A disciplined extension of the standard model}

The standard consumer model begins with coherent preferences, a budget
constraint, and purposeful choice. It is an exceptionally useful benchmark:
it tells us what a person would choose if they knew the relevant options and
could compare them without important limits on time, attention, memory, or
self-control.

\keyterm{Behavioural economics} asks what changes when those conditions are
only approximately true. It combines economic reasoning with evidence about
judgement and decision-making. The point is not to announce that people are
irrational. The point is to identify systematic patterns that matter for
prediction, welfare, or policy.

\begin{figure}[htbp]
  \centering
  \begin{tikzpicture}[font=\small]
    \node[font=\sffamily\bfseries,text=PPEBurgundy,anchor=east] at (-0.25,1.25)
      {Useful benchmark};
    \node[ppebox,minimum width=2.45cm,text width=2.15cm] (information) at (1.45,1.25)
      {Relevant\\information};
    \node[ppebox,minimum width=2.45cm,text width=2.15cm] (calculation) at (4.65,1.25)
      {Deliberation and\\comparison};
    \node[ppebox,minimum width=2.45cm,text width=2.15cm] (choice) at (7.85,1.25)
      {Chosen\\action};
    \draw[ppearrow] (information) -- (calculation);
    \draw[ppearrow] (calculation) -- (choice);

    \node[font=\sffamily\bfseries,text=PPENavarraRed,anchor=east] at (-0.25,-1.35)
      {Behavioural question};
    \node[ppebox,minimum width=2.45cm,text width=2.15cm] (limits) at (1.45,-1.35)
      {Limited time,\\attention, and memory};
    \node[ppebox,minimum width=2.45cm,text width=2.15cm] (shortcut) at (4.65,-1.35)
      {Heuristic, default,\\or self-control problem};
    \node[ppebox,minimum width=2.45cm,text width=2.15cm] (outcome) at (7.85,-1.35)
      {Predictable\\choice pattern};
    \draw[ppearrow] (limits) -- (shortcut);
    \draw[ppearrow] (shortcut) -- (outcome);
  \end{tikzpicture}
  \caption{Behavioural economics retains purposeful choice but investigates
  how limits on attention, information, and self-control shape it.}
  \label{fig:behavioural-benchmark}
\end{figure}
\FloatBarrier

\keyterm{Bounded rationality} means that people pursue objectives using limited
information and limited cognitive resources. A rule of thumb can be sensible
when a fully optimising calculation is costly. Behavioural analysis becomes
valuable when the same shortcut produces a predictable error across many
people or situations.

\begin{keyidea}
\textbf{A deviation is not automatically a mistake.}
Before calling a choice biased, ask what information, time, risk, social
meaning, or transaction cost the simplified model has omitted. The strongest
behavioural explanations identify a repeatable pattern and a plausible
decision process, rather than merely describing a surprising outcome.
\end{keyidea}

\section{Thinking, Fast and Slow}

Daniel Kahneman's phrase \emph{Thinking, Fast and Slow} offers a helpful
starting point for behavioural economics. \keyterm{System 1} refers to fast,
automatic, associative judgement: recognising a familiar face, completing a
common phrase, or reacting to an obvious hazard. \keyterm{System 2} refers to
slower, attention-demanding reflection: comparing contracts, checking an
algebraic calculation, or resisting an appealing first answer.

\begin{center}
\small
\renewcommand{\arraystretch}{1.18}
\begin{tabular}{@{}p{3.05cm}p{3.55cm}p{3.85cm}@{}}
  \toprule
   & System 1 & System 2\\
  \midrule
  Typical operation & Fast, automatic, associative, and low effort & Slower, reflective, and attention demanding\\
  Strength & Handles familiar routine tasks efficiently & Can check rules, compare evidence, and perform calculation\\
  Vulnerability & Can rely too heavily on a salient cue or first impression & Is limited; fatigue, time pressure, and distraction make sustained checking difficult\\
  \bottomrule
\end{tabular}
\end{center}

These labels are a teaching shorthand, not a literal map of two separate
brain systems. Nor is fast judgement simply irrational or slow reflection
automatically correct. A skilled doctor may recognise a familiar pattern
quickly; a careful calculation can still use the wrong data. The economic
point is that the context of a decision can affect how much attention is used
and which errors become likely.

\keyterm{Selective attention} is especially important. Focusing on a
headline price, a vivid feature, or one item in a long form can leave other
relevant information unnoticed. \keyterm{Metacognition}---reflecting on the
reliability of one's own reasoning---can help people decide when to pause,
seek a second estimate, or calculate explicitly. High-stakes, unfamiliar,
numerical, or counterintuitive choices are often good candidates for this
deliberate check.

\begin{exercisebox}
\textbf{Pause and check the first answer.}
\begin{enumerate}[label=(\alph*)]
  \item A bat and a ball cost \euro1.10 in total. The bat costs \euro1 more
  than the ball. The ball costs \textbf{\euro0.05}, not \euro0.10.
  \item A patch of lilies doubles in area each day and covers a lake on day
  48. It covers half the lake on \textbf{day 47}.
  \item If a runner passes the person in second place, the runner is now
  \textbf{second}, not first.
\end{enumerate}
The intuitive answers are attractive because they are quick, not because the
questions measure intelligence. The lesson is to recognise when a tempting
answer deserves a short System-2 check.
\end{exercisebox}

\section{Heuristics and judgement}

\keyterm{Heuristics} are mental shortcuts. They allow people to make fast
judgements without processing every available fact. Three are particularly
important in introductory behavioural economics.

\begin{center}
\small
\renewcommand{\arraystretch}{1.18}
\begin{tabular}{@{}p{2.75cm}p{3.35cm}p{4.50cm}@{}}
  \toprule
  Heuristic or bias & Basic idea & Illustrative risk\\
  \midrule
  Availability & Easily recalled examples receive excessive weight. & A vivid recent event may make a rare risk feel common.\\
  Anchoring & An initial number influences later judgement, even when it is not very informative. & The first advertised price can shape a buyer's estimate of a reasonable price.\\
  Representativeness & People judge likelihood by resemblance to a stereotype and may neglect base rates. & A persuasive story can dominate information about how frequent an event really is.\\
  \bottomrule
\end{tabular}
\end{center}

These shortcuts can be adaptive. If information is expensive and a quick
decision is needed, an available example or a familiar pattern may be useful.
The issue is whether the shortcut remains reliable in the environment at hand.
For a policy maker, a financial adviser, or a platform designer, predictable
errors raise the question of whether the choice environment should make
relevant information easier to notice.

\subsection{The main judgement shortcuts, with examples}

\textbf{Anchoring and adjustment.} An \keyterm{anchor} is an initial number
or starting point that shifts later judgement more than its information
warrants. A laptop marked as formerly \euro1,600 and now \euro1,200 can seem
more valuable than the same laptop displayed only at \euro1,200, even when the
crossed-out price is not a credible market price. A useful response is to
compare independent prices or calculate value from relevant features rather
than adjusting from the seller's first number.

\textbf{Availability.} The \keyterm{availability heuristic} gives too much
weight to examples that are easy to retrieve from memory. After vivid news
coverage of a rare accident, a traveller may overestimate its probability.
Easy recall can contain useful information, but recall also depends on
recency, emotion, and media coverage. A base rate or a well-defined comparison
class is often a better guide to frequency.

\textbf{Representativeness and base rates.} The
\keyterm{representativeness heuristic} judges likelihood by resemblance to a
typical case. Suppose a shy person is described as resembling a librarian.
That resemblance does not by itself make the person likely to be a librarian
if salespeople are vastly more numerous. The error is \keyterm{base-rate
neglect}: a stereotype crowds out information about prevalence.

\textbf{Regression to the mean.} \keyterm{Regression to the mean} is a
statistical tendency, not a bias: an unusually high or low noisy measurement
is, on average, followed by a measurement closer to the underlying average.
The judgement error is to assign a causal explanation to that ordinary return.
For example, a manager may conclude that praise worsens performance when an
exceptional worker performs less exceptionally the next week, or that criticism
works when an unusually poor worker subsequently improves. Before changing a
policy, compare similar workers who did and did not receive the intervention.

\subsection{Confidence, missing evidence, and persuasive stories}

\textbf{Overconfidence, optimism, and planning fallacy.}
\keyterm{Overconfidence} is excessive certainty about one's knowledge or
accuracy. \keyterm{Optimism bias} is an overly favourable expectation about
future outcomes. The \keyterm{planning fallacy} is a related tendency to
underestimate time and costs while overestimating benefits. A student team may
forecast that a research assignment will take four days and then need eight:
the practical remedy is to use data from comparable completed projects, not
only the team's best-case plan.

\textbf{WYSIATI.} What You See Is All There Is
(\keyterm{WYSIATI}) describes judging a situation mainly from the evidence
that is salient and available, without adequately asking what relevant
evidence is missing. A glossy company presentation may emphasise rising sales
while omitting debt, pending litigation, or the sales of competitors. The
right response is not paralysis; it is to ask what information would change
the conclusion and whether it has been checked. Rare, poorly understood
possibilities are sometimes called black swans, but that label is a warning
about model uncertainty, not a separate cognitive bias.

\textbf{Halo effects and narrative fallacies.} A \keyterm{halo effect} occurs
when one favourable or unfavourable impression influences an unrelated
judgement: treating an applicant as technically strong solely because they
attended a prestigious university, for example. A \keyterm{narrative fallacy}
is the temptation to construct a coherent causal story from sparse or noisy
events, such as explaining a one-day share-price movement with an untested
headline. Both invite a disciplined question: which evidence directly supports
the conclusion, and which evidence merely makes the story feel persuasive?

\section{Reference points, framing, and loss aversion}

Standard consumer theory ranks final bundles. Behavioural evidence often
suggests that people also evaluate changes relative to a
\keyterm{reference point}: the status quo, an expectation, a previous price,
or a socially salient entitlement. \keyterm{Framing} occurs when equivalent
descriptions make different reference points or features of a choice salient.

\keyterm{Loss aversion} is the tendency for a loss relative to a reference
point to have greater psychological weight than an equivalent gain. This
helps explain reluctance to sell an item one already owns, strong reactions to
a price increase, and the political difficulty of removing an established
benefit. It does not imply that every person always dislikes losses by the
same amount.

\begin{exercisebox}
\textbf{A lost ticket is not only a loss-aversion example.}
Imagine that a cinema ticket costs \euro10. One person loses a prepaid ticket
and hesitates to buy another; another loses a \euro10 note on the way to the
cinema but still buys a ticket. Both have lost \euro10 relative to their
starting resources. The different responses can reflect \keyterm{mental
accounting}: treating the ticket purchase and lost cash as entries in separate
mental budgets. Mental accounting can interact with loss aversion, but it is
more precise to distinguish the two mechanisms.
\end{exercisebox}

A common illustrative \keyterm{prospect-theory value function} is
\[
v(x)=
\begin{cases}
x^\alpha, & x\geq 0,\\
-\lambda(-x)^\alpha, & x<0,
\end{cases}
\qquad
0<\alpha<1,\quad \lambda>1,
\]
where \(x\) is a gain or loss relative to the reference point. The exponent
\(\alpha\) captures diminishing sensitivity, while \(\lambda\) captures the
extra weight on losses. Figure~\ref{fig:reference-dependent-value} shows the
shape, not a universal numerical estimate.

\begin{figure}[htbp]
  \centering
  \begin{tikzpicture}[x=0.95cm,y=0.78cm,font=\small]
    \draw[-{Latex[length=2.0mm]},PPEInk] (-5.0,0) -- (5.15,0)
      node[below right] {change relative to reference point};
    \draw[-{Latex[length=2.0mm]},PPEInk] (0,-4.35) -- (0,3.65)
      node[above left] {subjective value};
    \draw[dashed,PPEMuted] (0,-4.0) -- (0,3.3);
    \node[below,text=PPEMuted] at (0,-0.25) {reference point};
    \draw[PPENavarraRed,line width=1.15pt,smooth]
      plot coordinates {(-4.55,-4.00) (-3.50,-3.55) (-2.40,-2.95)
      (-1.35,-2.15) (-0.55,-1.22) (0,0) (0.65,0.70) (1.45,1.25)
      (2.65,1.78) (4.55,2.28)};
    \node[text=PPEBurgundy,align=center] at (-2.65,-1.15)
      {losses:\\steeper};
    \node[text=PPEBurgundy,align=center] at (2.85,2.75)
      {gains:\\diminishing sensitivity};
  \end{tikzpicture}
  \caption{An illustrative reference-dependent value function. The curve is
  steeper for losses than for gains and becomes flatter away from the
  reference point.}
  \label{fig:reference-dependent-value}
\end{figure}
\FloatBarrier

Framing changes neither the physical object nor the budget constraint by
itself. It changes attention and interpretation. For example, a policy
described as removing a benefit can trigger a different response from the
same policy described as ending a temporary discount. A careful analyst should
ask whether the descriptions are truly equivalent, what reference point each
one creates, and whose welfare criterion is being used. The same issue appears
when a medical treatment is described as having a \(95\%\) survival rate rather
than a \(5\%\) mortality rate. The figures are mathematically equivalent, but
the descriptions can make different aspects of risk salient.

\section{Present bias and self-control}

Many choices involve a conflict between an immediate cost and a delayed
benefit: studying, saving, exercising, or reducing carbon-intensive
consumption. \keyterm{Present bias} describes a tendency to give unusually
high weight to immediate consequences relative to future ones. A person may
sincerely plan in January to save more next month, yet repeatedly postpone the
same action when the moment arrives.

One simple representation is quasi-hyperbolic discounting:
\[
U_t=u(c_t)+\beta\sum_{k=1}^{\infty}\delta^k u(c_{t+k}),
\qquad
0<\beta\leq1,\quad 0<\delta<1.
\]
The usual discount factor \(\delta\) values future consequences less than
current ones. When \(\beta<1\), outcomes after the present receive an
additional discount. This is a model of present bias, not a claim that all
impatience is an error: people may have genuine reasons to prefer resources
now.

Present bias helps explain the appeal of commitment devices. Automatic saving,
calendar reminders, or a voluntary rule that makes an undesirable option
harder to choose can help a person carry out a plan they already endorse.
But an intervention should not assume that every delay is a self-control
failure. Flexibility can matter when incomes, health, or family obligations
change.

\begin{exercisebox}
\textbf{Why now can be special.}
Compare a choice between \euro10 today and \euro13 tomorrow with the same
one-day delay one year from now: \euro10 in one year versus \euro13 in one
year and one day. A person who takes \euro10 in the first comparison but
\euro13 in the second displays a changing trade-off when the immediate option
is present. That pattern is what present bias is designed to capture; ordinary
impatience alone need not produce it.
\end{exercisebox}

\section{Beyond individual bias: social preferences and influence}

Not every departure from the narrow model of material self-interest is a
cognitive bias. \keyterm{Social preferences} describe concern for other
people's payoffs, intentions, fairness, or relative position. In an ultimatum
game, one person proposes how to divide \euro100 and the other can accept or
reject; rejection gives both people zero. A purely money-maximising responder
would accept any positive amount, but very unequal offers are often rejected.
The observation may reflect fairness or reciprocal motives, rather than a
mistake in calculation.

Social influence is also distinct from anchoring. In a \keyterm{conformity}
experiment,
a person can publicly repeat an obviously incorrect group answer even without
an irrelevant numerical anchor. The course slide labelled Anchoring II
illustrates this different mechanism. For PPE, the distinction matters: a
policy may change behaviour by supplying information, by changing material
incentives, or by changing the social meaning of an action.

\section{Choice architecture and public policy}

\keyterm{Choice architecture} is the way options are presented and organised.
A \keyterm{default} is the option that applies if a person does nothing. A
\keyterm{nudge} changes choice architecture while preserving options and
without substantially changing monetary incentives. An opt-out pension plan,
a clear reminder, and a simple form can be nudges when people remain free to
choose otherwise.

\keyterm{Status quo bias} is the tendency to retain an existing option partly
because changing requires attention and effort, or because change feels like a
loss. A worker may remain in an unsuitable electricity tariff or postpone
changing a pension contribution even when a better option is available. This
helps explain why defaults can have large effects. It does not show that the
default expresses a person's true preference: a responsible design makes the
default transparent, useful for most affected people, and genuinely easy to
leave.

\begin{table}[htbp]
  \centering
  \small
  \renewcommand{\arraystretch}{1.18}
  \begin{tabular}{@{}p{2.40cm}p{3.60cm}p{4.60cm}@{}}
    \toprule
    Policy tool & Main mechanism & A central limitation or ethical question\\
    \midrule
    Information or reminder & Makes relevant facts or a planned action more noticeable & May be ignored and may not address an externality or affordability problem.\\
    Default or simplification & Reduces effort required for a chosen option & Who selected the default, and is opting out genuinely easy?\\
    Tax, subsidy, or price & Changes monetary incentives and can address external costs or benefits & Financial burdens and distributional effects require evaluation.\\
    Regulation or ban & Removes or constrains harmful options & Restricts choice and needs a strong justification and credible enforcement.\\
    \bottomrule
  \end{tabular}
  \caption{Behavioural tools complement, rather than replace, standard policy
  instruments.}
  \label{tab:behavioural-policy-tools}
\end{table}
\FloatBarrier

The distinction matters. A tax on pollution changes the price faced by a
decision-maker and can internalise an external cost. A reminder about a
polluting activity may improve attention but does not create the same
incentive. Likewise, a ban removes an option, while a default changes what
happens when people do not actively choose. Good policy identifies the actual
problem before selecting an instrument.

\begin{keyidea}
\textbf{A behavioural policy needs an ethical as well as an empirical case.}
\begin{enumerate}
  \item Is there credible evidence of a persistent decision problem?
  \item Does the intervention address its cause rather than merely its symptom?
  \item Is the choice architecture transparent and is opting out practical?
  \item Who bears the costs, who benefits, and could the design worsen inequality?
  \item Does the intervention have democratic authority and a realistic way to be evaluated?
\end{enumerate}
\end{keyidea}

Behavioural economics therefore enriches PPE rather than settling its political
or philosophical questions. It offers better descriptions of choice in some
settings, but it also forces us to ask who should design environments, what
counts as a person's own interest, and when influence becomes manipulation.

\section*{Chapter Review}

\subsection*{Key concepts and equations}
\begin{itemize}
  \item \textbf{Bounded rationality}: purposeful choice under limits on information, attention, time, or computation.
  \item \textbf{System 1 and System 2}: a teaching shorthand for fast, automatic judgement and slower, effortful reflection; not two literal brain modules.
  \item \textbf{Heuristic}: a mental shortcut that can save effort but can also create systematic error.
  \item \textbf{Anchoring, availability, and representativeness}: judgement patterns that can overweight an initial number, vivid examples, or resemblance to a stereotype; representativeness can lead to base-rate neglect.
  \item \textbf{Regression to the mean}: an ordinary statistical tendency that is often mistaken for a causal effect.
  \item \textbf{Overconfidence, optimism bias, and planning fallacy}: excessive certainty or favourable forecasts, including underestimation of time and cost.
  \item \textbf{WYSIATI}: forming a judgement from salient evidence without adequately asking what relevant evidence is missing.
  \item \textbf{Reference dependence}: evaluation of outcomes as gains or losses relative to a reference point.
  \item \textbf{Loss aversion and mental accounting}: losses receive greater weight than equal-sized gains relative to a reference point; mental accounts can separately frame equivalent losses.
  \item \textbf{Present bias}: extra weight placed on immediate consequences relative to later consequences.
  \item \textbf{Status quo bias}: inertia in retaining an existing option can amplify the effect of a default.
  \item \textbf{Nudge}: a choice-architecture intervention that preserves options and does not substantially alter monetary incentives.
\end{itemize}

\section*{Exercises}

\subsection*{Conceptual questions}
\begin{enumerate}[label=\arabic*.]
  \item Why is rational-choice analysis still a useful benchmark for behavioural economics?
  \item Identify the likely heuristic or bias in each case: a traveller overestimates a rare risk after vivid news coverage; a buyer's valuation is influenced by an initial list price; and a recruiter ignores the base rate when a candidate resembles a stereotype.
  \item Explain the relationship between framing, reference points, and loss aversion.
  \item Does evidence of a bias show that people are unintelligent? Explain using bounded rationality.
  \item Distinguish a nudge from a tax and a ban. Give one example of each.
\end{enumerate}

\subsection*{Graphical exercises}
\begin{enumerate}[resume,label=\arabic*.]
  \item Draw a reference-dependent value function that is steeper for losses than for gains. Mark the reference point, loss aversion, and diminishing sensitivity.
  \item Draw the behavioural decision pathway from limited attention or self-control to a heuristic, default, or commitment device, then to a choice. Give one concrete example for each stage.
\end{enumerate}

\subsection*{Analytical exercises}
\begin{enumerate}[resume,label=\arabic*.]
  \item In the simple value function \(v(x)=x\) for gains and \(v(x)=-2(-x)\) for losses, calculate \(v(\text{\euro10})\) and \(v(-\text{\euro10})\). How large must a gain be to offset the perceived loss of \euro10 in this example?
  \item A student faces an immediate study cost of \euro10 and a benefit of \euro15 next period. With \(\delta=0.95\), compare the choice under standard discounting with the choice under present bias \(\beta=0.60\). Calculate the net present value in each case.
  \item A pension plan has participation of \(35\%\) under active opt-in and \(75\%\) under automatic enrolment with opt-out. Calculate the change in percentage points. What evidence would be needed before attributing all of the difference to the default?
  \item A clinic wants to reduce missed appointments, which also prevent other patients from using scarce slots. Compare a reminder-and-easy-rescheduling policy with a financial no-show charge. State one advantage, one limitation, and one fairness question for each.
\end{enumerate}

\subsection*{Exam-style multiple-choice questions}
\begin{enumerate}[resume,label=\arabic*.]
  \item The availability heuristic involves: (a) using all information equally; (b) relying too heavily on easily recalled information; (c) ignoring memory; or (d) maximising expected utility exactly.
  \item Loss aversion means that: (a) gains and losses are valued equally; (b) a loss relative to a reference point often has more weight than an equal-sized gain; (c) people never take risks; or (d) prices cannot affect choice.
  \item Which policy is most clearly a nudge? (a) Automatically enrolling workers in a pension plan while making opt-out easy; (b) forbidding private saving; (c) taxing all savings; or (d) requiring everyone to buy the same pension product.
\end{enumerate}

\subsection*{PPE discussion question}

\noindent Are nudges compatible with individual autonomy? Evaluate a real or
plausible default policy using transparency, ease of opting out, evidence of a
decision problem, distributional effects, and democratic legitimacy.
